\section{Discussion}

The results of this study suggest that LLM-based molecular 
property prediction exhibits systematic sensitivity to structured 
semantic perturbations, with the direction and magnitude of 
observed effects depending on the task domain.

\subsection{Textual Augmentation as Semantic Conditioning}
A primary finding of this study is that the \textit{semantic 
orientation} induced by a prompt---the topical focus on specific 
molecular features---is a consistent driver of predictive 
shifts, often more so than the factual taxonomic correctness of the generated 
text. This was observed in the C4a (Structural) condition, 
where the LLM generated contextually plausible descriptions 
focused on structural features but failed to satisfy the strict 
Structural Phantom definition (i.e., referencing absent 
elements). Despite this taxonomic non-compliance, C4a 
consistently yielded the largest directional ROC-AUC 
improvements on BBBP and Tox21 while preserving baseline 
discriminability on BACE and yielding a marginal directional improvement over C0 on ESOL (RMSE: 2.045 vs.\ 2.051).

We frame this pattern as an effect of \textit{semantic conditioning}: the 
observation that factually non-compliant but topically 
structured text can systematically shift prediction behavior 
in a task-dependent direction. We emphasize that this 
characterization is descriptive and behavioral---it refers to 
the observed co-occurrence of structural topic focus and 
directional performance shifts, and does not constitute a 
mechanistic claim about the model's internal processing.

\subsection{Distributional Compression on Enzymatic and 
Solubility Tasks}
The consistent performance degradation observed on BACE and ESOL provides evidence that textual augmentation is systematically harmful in specific task regimes, though the effect is modulated by both vocabulary and molecular specificity. Notably, the enzymatic regime appears particularly sensitive to chemical-domain vocabulary presented in a generic, non-molecule-specific context---as evidenced by the significant degradation in C2 ($p_{\text{BH}}=0.011$) and factual C1 ($p_{\text{BH}}=0.011$). Conversely, pure non-chemical gibberish (C2b, $p_{\text{BH}}=0.217$) and molecule-specific structural descriptions (C4a, $p_{\text{BH}}=0.808$) did not induce statistically significant performance loss. This pattern suggests that the disruptive mechanism on enzymatic tasks is primarily driven by domain-relevant but generic contextual noise that displaces the SMILES representation, while specific structural anchoring (C4a) or completely irrelevant register (C2b) is less disruptive. This indicates that while fabrication leads to entropy collapse, distributional noise from generic chemical-rich text is equally disruptive to prediction ranking.

Quantitatively, this is reflected in the calibration metrics 
reported in Table~\ref{tab:calibration}. Under the C3 
condition on BACE, Brier Score increased from 0.244 to 0.438 
and ECE from 0.151 to 0.439---indicating that hallucination 
augmentation does not merely introduce ranking noise, but 
systematically degrades probability calibration. This 
behavioral pattern, which we term \textit{distributional 
compression}, consists of a collapse of the model's output 
distribution toward a narrow low-confidence peak near 
$p \approx 0.17$, regardless of true label. This is distinct 
from regression toward the BACE test set's class prior 
(60.5\% positive), indicating a genuine loss of 
discriminative signal rather than label-frequency anchoring.

One plausible account is that fabricated enzymatic or solubility claims introduce conflicting contextual signals that override the limited structural information conveyed by the SMILES input, collapsing the model's output toward a non-discriminative state. However, we note that this account remains speculative in the absence of mechanistic interpretability analysis. Crucially, the enzymatic regime is sensitive to any textual augmentation that shifts the model's attention away from SMILES-derived features. A notable exception is the prompt-constrained structural condition (C4a), which preserved baseline AUC (0.666 vs.\ 0.674) and substantially improved calibration (ECE: 0.069 vs.\ 0.151). This asymmetry suggests that the structural topic focus of C4a may anchor the model's attention to features that are complementary to SMILES parsing, whereas generic descriptor strings or fabricated claims introduce distributional noise that overwhelms the model's limited enzymatic reasoning capacity.

\subsection{Scaling Divergence and Representational States}
A critical observation in our scaling exploration (Section~\ref{sec:implementation_robustness}) is the divergence in behavior between Llama-3.1-8B and Llama-3.3-70B on the BBBP task. While the 8B model exhibited directional improvements with structural augmentation, the 70B model showed directional degradation when evaluated on the same textual augmentations. This reversal suggests that the ``utility'' of semantic augmentation is not an intrinsic property of the text itself, but is a function of the model's baseline representational state. If a larger model already captures the necessary structural information from the SMILES input, the addition of textual descriptions (especially those generated by a smaller 8B model) may introduce harmful interference rather than helpful conditioning. This underscores the necessity of task-specific and model-specific calibration when deploying text-augmented molecular reasoning systems.

\subsection{Generalization and Pre-training Artifacts}
The contrast between BBBP/Tox21 and BACE/ESOL suggests that 
task sensitivity may correlate with the prevalence of 
structural reasoning vs.\ specific mechanistic retrieval 
in the pre-training distribution. General physicochemical 
properties---such as blood-brain barrier permeability and 
broad toxicity---are more broadly represented in the 
scientific literature in association with diverse molecular 
scaffolds. In these domains, structural topic focus (C4a) 
may prime the model's contextual inference toward relevant 
sub-graph patterns, producing directional ROC-AUC 
improvements even without factual accuracy. Specific enzymatic mechanisms and 
quantitative solubility, by contrast, may be less densely 
covered in general-domain pre-training corpora, making 
predictions in these regimes more susceptible to disruption 
by fabricated claims.

We note that this interpretation is speculative and cannot 
be verified without targeted analysis of training data 
composition or model internals, which is beyond the scope 
of this exploratory study.

The observed cross-task consistency also suggests that 
semantic perturbation experiments of this kind may serve 
as a lightweight behavioral probe for characterizing 
task-specific inference sensitivity in LLMs---one that 
does not require access to model weights or internal 
representations, and can be applied to black-box API-based 
systems. We make no stronger claim about the mechanistic 
basis of this sensitivity.

\subsection{Methodological Implications}
Our findings carry several practical implications for 
LLM-augmented cheminformatics. First, the observed 
directional trend of structural topic-focus prompting (C4a) over 
factual RDKit descriptors (C1) on physicochemical tasks 
suggests that topically oriented textual context may 
sometimes provide a more effective inference substrate for 
general-purpose LLMs than numerical descriptor strings, 
which may not integrate naturally into language-based 
reasoning chains. Second, the severe calibration failure 
on BACE---reflected in ECE values approaching 0.44---
underscores the risk of deploying LLM-augmented 
predictions in enzymatic domains without explicit 
uncertainty quantification. Future work should examine 
whether calibration-aware prompting strategies or 
post-hoc recalibration can mitigate this effect.
